% !TeX program = xelatex
% UTF-8 编码。使用完整 TeX Live 或 Overleaf，选择 XeLaTeX 编译。
% 使用 TeX Live 自带的 Fandol 中文字体，无需额外字体或参考文献文件。
\documentclass[UTF8,a4paper,zihao=-4,fontset=fandol]{ctexart}
\usepackage[margin=2.5cm,includehead=false,headheight=0pt,headsep=0pt]{geometry}
\usepackage{amsmath,amssymb}
\usepackage{enumitem}
\usepackage[hidelinks]{hyperref}
\usepackage{booktabs}
\usepackage{tabularx}
\setlist[enumerate]{leftmargin=2em,itemsep=0.25em,topsep=0.3em}
\setlength{\parskip}{0.25em}
\linespread{1.08}
\ctexset{section={format=\large\bfseries,beforeskip=1.2ex,afterskip=0.6ex}}
\pagestyle{plain}
\title{\vspace{-1.5cm}\textbf{生成式推荐中的兴趣与决策倾向}}
\author{}
\date{}

\begin{document}
\maketitle
\vspace{-1.6cm}

\section{问题}

生成式推荐中，SID 可视为物品在不同粒度上的语义表征。根据用户历史生成 SID，可以理解为逐步确定用户可能感兴趣的范围，最终得到完整的 SID 对应具体物品。训练和评估通常以实际发生的下一次交互作为目标。但是唯一的 next item 标签不等于唯一的兴趣：

多个物品都可能符合用户当时的兴趣，实际交互只是用户的最终选择。例如，几件用途相近的护肤品都可能符合需求，但用户最终选择哪一件，还可能受到相对价格、质量、流行度等因素的影响。

因此，我觉得其中一个可以进一步研究的问题是：通过分析“在符合用户兴趣的物品范围内，哪些决策倾向影响了用户的最终选择”，来提高模型在 next item prediction 中的精准度。

兴趣用于确定可能感兴趣的物品范围，决策倾向用于解释和预测这一范围内的取舍。将两者区分，研究仅靠兴趣匹配尚未充分解释的选择差异。

\section{决策倾向}

决策倾向指用户在可比物品之间，对不同选择表现出的取舍规律。例如，用户在护肤品、美妆和厨具中，都经常选择同类中相对便宜的产品，那么可能存在跨越不同兴趣的低价倾向。这里的研究对象是选择时对价格的取舍，而非偏好护肤、美妆或厨具这类品类兴趣。

同一用户可能同时考虑多个因素，其重要程度也可能随场景变化。是否存在跨兴趣、跨时间的稳定倾向，应通过历史选择和后续预测检验，不预先为每位用户指定一个固定的主导因素。

\section{初步思路}
\subsection{}
既然 SID 能量化物品的语义相似度，那生成的 SID 本质上映射的是用户兴趣偏好，而最终的 next item prediction 其实是决策层面的，应该由用户的决策主导因素来决定（也就是：SID 管兴趣，决策主导因素管最终选择，两者分开）。

\subsection{}
把这些决策主导因素作为 Transformer 中 QKV 里 Q 的一部分（K 也许也可以），也就是在 transformer 里用上这些主导因素。

\section{评估}

评估仍以 next item prediction 为目标，采用 Recall@$K$ 和 NDCG@$K$。

\section{验证实验}

\subsection{实验设置}
通过以下实验验证决策因素的存在及重要性：

传统的生成式推荐，如 TIGER，当存在 SID 冲突时会采用各种方式避免，让每个 item 对应的 SID 组合 unique。在我们的既定假设中，SID 其实反映的是用户的兴趣倾向，故冲突的 SID 组合可以认为物品属于同一兴趣大类。比如：

洗发水A 与 洗发水B 通过 RQ-VAE 生成出的 SID 组合均为 (1,2,3)。那在TIGER中，会通过再随机分配一位来避免冲突，如 A(1,2,3,0)，B(1,2,3,1)。

而我们不进行冲突避免的操作，(1,2,3) 表示 A、B都同属洗发水这一兴趣品类。

假定用户 u 的 ground truth 是 A，而模型生成的是 B，由于二者对应的 SID 组合一致（均为(1,2,3)），故认为模型的兴趣偏好判断正确。

但由于实际是 A，而非 B。故，如果此时模型各结果指标相较于有冲突避免时有所提升，则表明最后的错误只是由于我们对用户 u 在做最终决策偏好时的决策因素的判断错误所造成。

这也就论证了决策因素是 GR 在 next item predition 任务中很重要的一环。

\subsection{评估指标及计算}

\paragraph{符号定义}
设用户 $u$ 的测试目标物品为 $y_u$，
模型输出的物品排序为
$P_u=(p_{u,1},p_{u,2},\ldots)$。
$s(i)$ 表示物品 $i$ 的前 $D$ 层 SID，
$c(i)$ 表示追加消冲突层后的唯一编码。
共享同一 SID 的物品构成集合
\[
    \mathcal{C}(i)=\{j:s(j)=s(i)\}.
\]
其中，$|\mathcal{C}(i)|>1$ 表示该 SID 存在编码冲突。

\paragraph{三种排名}
对于 \texttt{sid\_nc}，定义：
\begin{align}
    r_u^{\mathrm{item}}
    &= \min\{j:c(p_{u,j})=c(y_u)\},\\
    r_u^{\mathrm{sid\text{-}item}}
    &= \min\{j:s(p_{u,j})=s(y_u)\},\\
    r_u^{\mathrm{sid\text{-}unique}}
    &= \operatorname{rank}\!\left(
        s(y_u),
        \operatorname{UniqueInOrder}
        [s(p_{u,1}),s(p_{u,2}),\ldots]
    \right).
\end{align}
未匹配时，排名记为 $+\infty$。
$\operatorname{UniqueInOrder}$ 表示按首次出现顺序去重。

三种排名分别对应：
\begin{itemize}
    \item \texttt{item}：精确目标物品的排名；
    \item \texttt{sid\_at\_item\_rank}：
    首个与目标共享 SID 的物品位置；
    \item \texttt{sid\_unique}：
    将预测物品映射为 SID 并去重后，目标 SID 的排名。
\end{itemize}

对于 \texttt{sid}，模型直接预测 SID 类，
以目标 SID 在预测列表中的位置计算
\texttt{sid/*} 指标。
其排名单位与 \texttt{sid\_unique} 相同，
但训练表示和生成过程不同，结果不保证一致。

\paragraph{Recall 与 NDCG}
对于任意一种排名 $r_u$，
单个测试样本的指标为：
\begin{align}
    \operatorname{Recall@K}(u)
    &= \operatorname{Hit@K}(u)
     = \mathbb{I}[r_u\leq K],\\
    \operatorname{NDCG@K}(u)
    &=
    \begin{cases}
        \dfrac{1}{\log_2(r_u+1)}, & r_u\leq K,\\
        0, & \text{否则}.
    \end{cases}
\end{align}
其中，$K$ 为评估截断位置，
$\mathbb{I}[\cdot]$ 为指示函数。
每个用户只有一个测试目标，因此 Recall 与 Hit 相等；
NDCG 进一步反映命中位置，目标排名越靠前，得分越高。

\paragraph{SID 命中但物品未命中}
定义：
\[
    \Delta_u@K
    =
    \mathbb{I}[r_u^{\mathrm{sid\text{-}item}}\leq K]
    -
    \mathbb{I}[r_u^{\mathrm{item}}\leq K].
\]
该指标对应 \texttt{sid\_hit\_item\_miss/recall@K}。
当目标 SID 已在前 $K$ 个物品中出现，
但精确目标物品未出现时，其值为 1；否则为 0。
其平均值表示此类情况占测试样本的比例。

\paragraph{计算示例}
设 $A_1,A_2$ 共享 SID $S_A$，
$B_1,B_2$ 共享 SID $S_B$，
$C_1$ 的 SID 为 $S_C$。
真实目标为 $B_1$，预测排序为：
\[
    P_u=(A_1,A_2,B_2,C_1,B_1).
\]
对应 SID 序列及其去重结果为：
\[
    (S_A,S_A,S_B,S_C,S_B)
    \;\longrightarrow\;
    (S_A,S_B,S_C).
\]
因此：
\[
    r_u^{\mathrm{item}}=5,\qquad
    r_u^{\mathrm{sid\text{-}item}}=3,\qquad
    r_u^{\mathrm{sid\text{-}unique}}=2.
\]

\begin{center}
\begin{tabular}{lccc}
    \hline
    指标 & \texttt{item}
         & \texttt{sid\_at\_item\_rank}
         & \texttt{sid\_unique}\\
    \hline
    Recall@3 & 0 & 1 & 1\\
    NDCG@3   & 0 & 0.5 & 0.63093\\
    Recall@5 & 1 & 1 & 1\\
    NDCG@5   & 0.38685 & 0.5 & 0.63093\\
    \hline
\end{tabular}
\end{center}

在此例中，$\Delta_u@3=1$：
同 SID 的 $B_2$ 已进入前 3 位，
但精确目标 $B_1$ 尚未进入。

\paragraph{分组与汇总}
设测试用户集合为 $\mathcal{U}$，按目标 SID 类大小划分：
\begin{align}
    \mathcal{U}_{\mathrm{all}}
    &= \mathcal{U},\\
    \mathcal{U}_{\mathrm{collision}}
    &= \{u:|\mathcal{C}(y_u)|>1\},\\
    \mathcal{U}_{\mathrm{singleton}}
    &= \{u:|\mathcal{C}(y_u)|=1\}.
\end{align}
SID 类大小依据整个物品映射计算。
模型仅预测一遍，每个样本的指标同时计入
\texttt{all} 和对应子组。
对任意分组 $\mathcal{V}$，最终指标为：
\[
    \overline{M}_{\mathcal{V}}
    =
    \frac{1}{|\mathcal{V}|}
    \sum_{u\in\mathcal{V}}M_u.
\]
因此，\texttt{all} 表示整体表现，
\texttt{collision} 表示目标存在 SID 冲突时的表现，
\texttt{singleton} 表示目标 SID 唯一时的表现。
统计单位为测试样本（用户），而非不同物品。

\end{document}